Professional Mathematics Teacher.
Solve a specific LaTeX exercise involving a function $f(x) = 2x^3 - 3ax^2 + 5$ ($a \ne 0$).

        *   Wrap solution in `\begin{thema}{Γ} ... \end{thema}`.
        *   Questions in `\begin{erwthma} ... \end{erwthma}`.
        *   Use `\item` for sub-questions.
        *   `$ ... $` for inline math, `\[ ... \]` for display math.
        *   Use `\begin{align*}` for multi-line calculations.
        *   Use `\eng{...}` for English terms.
        *   Use `\textbf{...}` for emphasis.
        *   Strict restrictions: No indefinite integrals, no 2nd derivative test for extrema, no differential equations, no MVT for integrals, no differentiation of integrals with variable limits.
        *   Formatting identifiers: `%=========== ΛΥΣΗ - Xο ΘΕΜΑ Γ ===========` and `%=========== ΤΕΛΟΣ ΛΥΣΗΣ - Χο ΘΕΜΑ Γ ===========`.
        *   No preamble, no comments inside solution, no repetition of the exercise text.

    *   $f(x) = 2x^3 - 3ax^2 + 5$ ($a \neq 0$).
    *   $f'(x) = 6x^2 - 6ax = 6x(x - a)$.

    *   **Part 1: Positions of local extrema.**
        *   $f'(x) = 0 \Leftrightarrow 6x(x - a) = 0 \Leftrightarrow x = 0$ or $x = a$.
        *   Since $a \neq 0$, these are two distinct roots.
        *   The positions are $x_1 = 0$ and $x_2 = a$.

    *   **Part 2: Locus of $M(a, f(a))$.**
        *   $M$ has coordinates $(x, y) = (a, f(a))$.
        *   $x = a$.
        *   $y = f(a) = 2a^3 - 3a(a^2) + 5 = 2a^3 - 3a^3 + 5 = 5 - a^3$.
        *   Substitute $a = x$ into $y = 5 - a^3$:
        *   $y = 5 - x^3$.

    *   **Part 3: Inflection point at $x=1$.**
        *   $f'(x) = 6x^2 - 6ax$.
        *   $f''(x) = 12x - 6a$.
        *   A point of inflection occurs where $f''(x) = 0$ and $f''(x)$ changes sign.
        *   For $x=1$ to be a point of inflection: $f''(1) = 0 \Rightarrow 12(1) - 6a = 0 \Rightarrow 6a = 12 \Rightarrow a = 2$.

    *   **Part 4: Limit for $a=2$.**
        *   Given $a=2$.
        *   $f(x) = 2x^3 - 6x^2 + 5$.
        *   $f'(x) = 6x^2 - 12x$.
        *   $f(x) - 5 = 2x^3 - 6x^2$.
        *   Limit: $\lim_{x \to 0} \frac{2x^3 - 6x^2}{6x^2 - 12x} = \lim_{x \to 0} \frac{2x^2(x - 3)}{6x(x - 2)} = \lim_{x \to 0} \frac{x(x - 3)}{3(x - 2)}$.
        *   Evaluation: $\frac{0(-3)}{3(-2)} = \frac{0}{-6} = 0$.

    *   Check against Guidelines:
        *   `\begin{thema}{Γ}` $\checkmark$
        *   `\begin{erwthma}` $\checkmark$
        *   `\item` $\checkmark$
        *   `\[ ... \]` and `\begin{align*}` $\checkmark$
        *   `\eng{...}` for English terms (e.g., $\eng{L'Hôpital}$ if used, though not needed here).
        *   `\textbf{...}` for key terms (e.g., $\textbf{τοπικά ακρότατα}$, $\textbf{σημείο καμπής}$).
        *   No prohibited concepts used.
%=========== ΛΥΣΗ - 3ο ΘΕΜΑ Γ ===========
\begin{thema}{Γ}
\begin{erwthma}
\item Η συνάρτηση $f$ είναι πολυωνυμική, άρα είναι παραγωγίσιμη στο $\mathbb{R}$. Η παράγωγός της είναι:
\[ f'(x) = 6x^2 - 6ax = 6x(x-a) \]
Για τον προσδιορισμό των \textbf{τοπικών ακροτάτων}, λύνουμε την εξίσωση $f'(x) = 0$:
\[ 6x(x-a) = 0 \Leftrightarrow x = 0 \text{ ή } x = a \]
Επειδή $a \neq 0$, οι δύο ρίζες είναι διακριτές. Επομένως, οι θέσεις των δύο τοπικών ακροτάτων είναι οι $x_1 = 0$ και $x_2 = a$.

\item Το σημείο $M(a, f(a))$ έχει συντεταγμένες:
\[ x = a \]
\[ y = f(a) = 2a^3 - 3a(a^2) + 5 = 2a^3 - 3a^3 + 5 = 5 - a^3 \]
Από την πρώτη σχέση έχουμε $a = x$. Αντικαθιστώντας στη δεύτερη σχέση, προκύπτει:
\[ y = 5 - x^3 \]
Άρα, καθώς το $a$ μεταβάλλεται, το σημείο $M$ κινείται πάνω στην καμπύλη με εξίσωση $y = 5 - x^3$.

\item Η δεύτερη παράγωγος της συνάρτησης $f$ είναι:
\[ f''(x) = 12x - 6a \]
Για να παρουσιάζει η γραφική παράσταση της $f$ \textbf{σημείο καμπής} στο $x=1$, πρέπει να ισχύει $f''(1) = 0$:
\[ 12(1) - 6a = 0 \Leftrightarrow 6a = 12 \Leftrightarrow a = 2 \]
Επειδή η $f''(x)$ είναι γραμμική συνάρτηση με $f''(1)=0$, αλλάζει πρόσημο στο $x=1$, οπότε για $a=2$ το σημείο με τετμημένη $x=1$ είναι πράγματι σημείο καμπής.

\item Για $a=2$, η συνάρτηση και η παράγωγός της είναι:
\[ f(x) = 2x^3 - 6x^2 + 5 \quad \text{και} \quad f'(x) = 6x^2 - 12x \]
Ο υπολογισμός του ορίου είναι:
\begin{align*}
\lim_{x \to 0} \frac{f(x) - 5}{f'(x)} &= \lim_{x \to 0} \frac{2x^3 - 6x^2}{6x^2 - 12x} \\
&= \lim_{x \to 0} \frac{2x^2(x - 3)}{6x(x - 2)} \\
&= \lim_{x \to 0} \frac{x(x - 3)}{3(x - 2)} \\
&= \frac{0 \cdot (-3)}{3 \cdot (-2)} = \frac{0}{-6} = 0
\end{align*}
\end{erwthma}
\end{thema}
%=========== ΤΕΛΟΣ ΛΥΣΗΣ - 3ο ΘΕΜΑ Γ ===========